\originalsection{作业2}\label{sec:ps02}
原件3页, 数学题面2页.
\begin{exercise}\label{ps02:1}
原题1. 设 $B$ 为 Banach 空间. (a) 若 $T\in\B(B,B)$ 且 $\norm{I-T}<1$, 证明 $T$ 可逆, 且 $\sum_{n=0}^\infty(I-T)^n$ 在算子范数中收敛于 $T^{-1}$. (b) 证明可逆算子的集合在 $\B(B,B)$ 中开.
\end{exercise}
\begin{proof}[AI 独立解答]
(a) 令 $A=I-T$, 几何级数界 $\sum\norm{A^n}\le\sum\norm A^n<\infty$ 与算子空间完备性保证和 $S$. 有 $T\sum_{n=0}^NA^n=(\sum_{n=0}^NA^n)T=I-A^{N+1}$; 取范数极限得 $TS=ST=I$. 这是定理\ref{thm:neumann}的证明, 此处只依赖算子空间完备性, 不依赖谱理论.

(b) 若 $T$ 可逆且 $\norm{T-S}<\norm{T^{-1}}^{-1}$, 则 $S=T[I-T^{-1}(T-S)]$, 方括号由(a)可逆, 从而 $S$ 可逆. 零空间的唯一算子可逆, 其算子空间本身也是开集, 另行满足结论.
\end{proof}
\begin{exercise}\label{ps02:2}
原题2. 设 $V$ 为赋范空间, $W\subsetneq V$ 为闭子空间. (a) 证明 $\norm{v+W}=\inf_{w\in W}\norm{v+w}$ 是 $V/W$ 上的范数. (b) 对每个 $\eps>0$, 证明存在 $v\in V$ 使 $\norm v=1$ 且 $\norm{v+W}\ge1-\eps$. 原提示: 对 $u\notin W$ 选接近距离下确界的代表元, 再归一化.
\end{exercise}
\begin{proof}[AI 独立解答]
(a) 改换代表元不改变被取下确界的陪集. 齐次性由 $W$ 为子空间得到; 对两个陪集各取误差小于 $\delta$ 的代表元, 用三角不等式再令 $\delta\downarrow0$, 得商范数的三角不等式. 若商范数为零, 则存在 $w_n\in W$ 使 $v+w_n\to0$, 闭性给出 $v\in W$; 反向显然. 此即定理\ref{thm:quotient}的范数部分.

(b) $\eps\ge1$ 时任取一个单位向量即可. $0<\eps<1$ 时选 $u\notin W$, 令 $d=\norm{u+W}>0$, 并选 $w\in W$ 使 $d\le\norm{u+w}<(1+\eps)d$. 则 $v=(u+w)/\norm{u+w}$ 为单位向量, 且 $\norm{v+W}=d/\norm{u+w}>1/(1+\eps)\ge1-\eps$.
\end{proof}
\begin{exercise}\label{ps02:3}
原题3. 若 $V$ 为 Banach 空间, $W\subsetneq V$ 为闭子空间, 证明上述商空间为 Banach 空间. 原提示: 对绝对可和的陪集级数选 $\norm{v_n+w_n}\le\norm{v_n+W}+2^{-n}$, 利用 $V$ 内的绝对级数收敛.
\end{exercise}
\begin{proof}[AI 独立解答]
按提示选代表元, 则 $\sum\norm{v_n+w_n}<\infty$. 定理\ref{thm:series}给出 $v=\sum(v_n+w_n)\in V$, 且
\[
\left\|v+W-\sum_{n=1}^N(v_n+W)\right\|\le\left\|v-\sum_{n=1}^N(v_n+w_n)\right\|\longrightarrow0.
\]
所以商空间中每个绝对级数收敛, 再用同一定理的逆向证明得完备. 这是定理\ref{thm:quotient}的完备性结论的另一种证明.
\end{proof}
\begin{exercise}\label{ps02:4}
原题4. 设 $V,W$ 为 Banach 空间, $T\in\B(V,W)$. (a) 证明 $\ker T$ 是闭子空间. (b) 称双射线性 $S:V_1\to V_2$ 为同构, 若 $S$ 及 $S^{-1}$ 都有界. 证明 $V/\ker T$ 与 $\ran T$ 同构当且仅当 $\ran T$ 闭. 原提示: 考察 $S(v+\ker T)=Tv$. 原件同构定义的箭头误写 $V_1\to V_1$, 此处按其定义域、值域及后续公式校正为 $V_1\to V_2$.
\end{exercise}
\begin{proof}[AI 独立解答]
(a) 线性性保证核为子空间, 连续性保证它是闭集 $\{0\}$ 的原像.

(b) $S$ 与代表元无关, 是线性双射, 且 $\norm{Tv}\le\norm T\norm{v+w}$ 对每个 $w\in\ker T$ 成立, 因而 $\norm S\le\norm T$. 若值域闭, 它与 $V/\ker T$ 都完备, 有界逆定理\ref{cor:inverse}给出 $S^{-1}$ 有界. 反之, 即使假设的是任何一个有界双向同构, 它也把完备的商空间映为完备空间: 值域的 Cauchy 列在逆像中 Cauchy, 再映回极限. 完备子空间在 $W$ 内闭, 故值域闭. 这不把“代数同构”误当成“拓扑同构”.
\end{proof}
\begin{exercise}\label{ps02:5}
原题5. 令 $W=\{a:\sum_{k\ge1}k|a_k|<\infty\}$, 配备 $\ell^1$ 范数. (a) 证明 $W$ 是 $\ell^1$ 的真稠密子空间, 从而不完备. (b) 定义 $(Ta)_k=ka_k$, 证明 $T:W\to\ell^1$ 的图闭, 但 $T$ 无界. (c) 证明 $S=T^{-1}:\ell^1\to W$ 有界且满射, 但不开. 原提示: 用有限截断证明稠密.
\end{exercise}
\begin{proof}[AI 独立解答]
(a) $c_{00}\subset W$, 截断在 $\ell^1$ 中逼近每个元素. 数列 $(k^{-2})$ 属于 $\ell^1$ 而不属于 $W$, 故 $W$ 真且不闭; 若完备就会在 $\ell^1$ 内闭, 矛盾.

(b) 若 $a^{(n)}\to a$ 于 $W$ 的所给范数, $Ta^{(n)}\to b$ 于 $\ell^1$, 坐标收敛给出 $b_k=ka_k$. 因 $b\in\ell^1$, $a\in W$ 且 $Ta=b$, 故图闭. 但 $\norm{Te_k}_1=k$, $\norm{e_k}_1=1$, 所以无界.

(c) $Sb=(b_k/k)_k$, 有 $\sum k|(Sb)_k|=\norm b_1$, $\norm{Sb}_1\le\norm b_1$, 故良定义、有界且与 $T$ 互逆. 若 $S$ 开, $S$ 的单位开球像应包含 $W$ 中某个半径 $\delta$ 的零点邻域. 取 $a=(\delta/2)e_k$ 且 $k>2/\delta$, 则 $\norm a_1<\delta$ 而 $\norm{Ta}_1>1$, 所以 $a$ 不在该像中, 矛盾. 这具体说明闭图定理不能去掉定义域完备性, 开映射定理不能去掉值域完备性.
\end{proof}
